\problemname{Proteinsyntes}
Latbo är väl känt att vara sveriges lataste stad. Björn har haft som mål att vara latast i stan
ända sedan han var liten. Igår kom han hela vägen till final på årets lathetstävling. Han var helt redo
på att vinna, men blev totalkrossad när hans motståndare var för lat för att ens dyka upp.

Detta var droppen för Björn och han har nu bestämt sig att använda sina medicinkunskaper för att
utveckla ett botemedel mot lathet. För att göra detta behöver han syntetisera det magiska proteinet $T$.
Han faller dock snabbt in i gamla vanor och försöker syntetisera det med hjälp av så få andra proteiner
som möjligt. 

Eftersom han är för lat för att bestämma hur han ska göra detta själv ber han nu dig om hjälp att 
hitta bästa sättet att konstruera protein $T$. Protein $T$ kan representeras som en sträng av längd $N$.
Björn känner till hur han kan producera $M$ sorters proteiner och vill minimera antalet protein han
behöver tillverka. När rätt mängd proteiner har tillverkats är det lätt att sätta ihop dessa till protein $T$. 

\section*{Indata}
Den första raden i indatan innehåller heltalet $N$ ($1 \le N \le 10^5$), längden på protein $T$.

Nästa rad innehåller en sträng av längd $N$, protein $T$.

På raden därefter följer $M$ ($1 \le M \le 10^6$), antalet protein Björn kan tillverka.

Därefter följer $M$ rader, som vardera innehåller en sträng. Denna sträng representerar ett protein han
kan tillverka. Det är garanterat att den totala längden av de $M$ strängarna inte överstiger $10^6$.

Alla strängar i indatan består av bokstäver \texttt{a}-\texttt{z}. Det är garanterat att Björn kan konstruera $T$ med hjälp av de $M$ strängar som Björn kan tillverka.

\section*{Utdata}
Skriv ut ett heltal: det minsta antalet protein som Björn behöver producera för att skapa protein $T$.


\section*{Poängsättning}
Din lösning kommer att testas på en mängd testfallsgrupper.
För att få poäng för en grupp så måste du klara alla testfall i gruppen.

\noindent
\begin{tabular}{| l | l | p{12cm} |}
  \hline
  \textbf{Grupp} & \textbf{Poäng} & \textbf{Gränser} \\ \hline
  $1$    & $10$       & $N \leq 20, M \leq 5$ \\ \hline
  $2$    & $20$       & $N, M \leq 100$ \\ \hline
  $3$    & $20$       & $M \leq 100$ \\ \hline
  $4$    & $25$       & $N \leq 30,000$ \\ \hline
  $5$    & $25$       & Inga ytterligare begränsningar. \\ \hline
\end{tabular}

\section*{Förklaring av exempelfall 1}
Genom att skapa proteinet \texttt{abc} två gånger kan han sätta ihop de och skapa protein $T$.
